% !TeX program = xelatex
% Compile twice with XeLaTeX. All illustrations are native vector TikZ drawings.
\documentclass[12pt,a4paper,oneside]{report}
\XeTeXgenerateactualtext=1
\usepackage[margin=23mm,top=24mm,bottom=24mm,headheight=27pt]{geometry}
\usepackage{amsmath,amssymb,mathtools,bm}
\usepackage{booktabs,array,longtable,tabularx}
\usepackage[table]{xcolor}
\usepackage{tikz}
\usetikzlibrary{arrows.meta,calc,patterns}
\usepackage[most]{tcolorbox}
\usepackage{titlesec,needspace}
\usepackage{hyperref}
\usepackage{bookmark}
\usepackage{xepersian}
\settextfont[Scale=1.06,BoldFont=Amiri-Bold.ttf]{Amiri-Regular.ttf}
\setlatintextfont[BoldFont=texgyretermes-bold.otf,ItalicFont=texgyretermes-italic.otf,BoldItalicFont=texgyretermes-bolditalic.otf]{texgyretermes-regular.otf}
\definecolor{navy}{HTML}{163747}
\definecolor{teal}{HTML}{127C80}
\definecolor{gold}{HTML}{B18542}
\definecolor{paper}{HTML}{F3F7F7}
\definecolor{muted}{HTML}{5E7178}
\definecolor{alert}{HTML}{A45538}
\hypersetup{unicode=true,colorlinks=true,linkcolor=navy,urlcolor=teal,
 pdftitle={روش‌های آماری ـ جزوهٔ بازنویسی‌شده},pdfauthor={علیرضا جهانبخش و احسان شهبازی},pdfsubject={احتمال، برآورد نقطه‌ای، فاصلهٔ اطمینان و آزمون فرض}}
\setlength{\parindent}{0pt}
\setlength{\parskip}{5pt plus 1pt}
\linespread{1.12}
\raggedbottom
\setcounter{tocdepth}{1}
\setcounter{secnumdepth}{2}
\allowdisplaybreaks[2]
\emergencystretch=1em
% Page styles use explicit boxes to keep bidi direction groups balanced.
\makeatletter
\def\ps@lecture{%
  \def\@oddhead{\vbox{\hbox to\textwidth{\small
    \rl{روش‌های آماری}\hfil\rl{\leftmark}}%
    \kern4pt\hrule height .35pt}}%
  \let\@evenhead\@oddhead
  \def\@oddfoot{\hfil\thepage\hfil}%
  \let\@evenfoot\@oddfoot}
\makeatother
\pagestyle{lecture}
\renewcommand{\chaptermark}[1]{\markboth{#1}{}}
\titleformat{\chapter}[display]{\color{navy}\bfseries\Huge}{\large فصل \thechapter}{8pt}{}[\vspace{4pt}\color{gold}\titlerule]
\titlespacing*{\chapter}{0pt}{0pt}{18pt}
\titleformat{\section}{\Large\bfseries\color{navy}}{\thesection}{.6em}{}
\titleformat{\subsection}{\large\bfseries\color{teal}}{\thesubsection}{.6em}{}
\tcbset{enhanced,breakable,boxrule=.45pt,arc=1.5mm,
 left=3mm,right=3mm,top=2mm,bottom=2mm,before skip=9pt,after skip=9pt,
 fonttitle=\bfseries,coltitle=white,colback=paper,colframe=teal}
\newtcolorbox{definition}[1]{breakable=false,title={تعریف | #1},colframe=navy,colback=navy!3!white}
\newtcolorbox{result}[1]{breakable=false,title={نتیجهٔ کلیدی | #1},colframe=teal,colback=teal!4!white}
\newtcolorbox{important}[1]{breakable=false,title={دقت کنید | #1},colframe=gold,colback=gold!6!white}
\newcounter{ex}[chapter]
\newenvironment{example}[1]{\refstepcounter{ex}\begin{tcolorbox}[breakable=false,title={مثال \thechapter.\arabic{ex} | #1},colframe=muted,colback=white]}{\end{tcolorbox}}
\newcommand{\edit}[1]{\par\begingroup\small\color{alert}\textbf{تکمیل یا اصلاح علمی.} #1\par\endgroup}
\newcommand{\source}[1]{\par{\small\color{muted} مبنا: صفحه‌های #1 فایل دست‌نویس.}\par}
\newcommand{\E}{\mathbb E}
\newcommand{\Var}{\operatorname{Var}}
\newcommand{\Cov}{\operatorname{Cov}}
\newcommand{\MSE}{\operatorname{MSE}}
\newcommand{\Bias}{\operatorname{Bias}}
\newcommand{\SE}{\operatorname{SE}}
\newcommand{\iid}{\overset{\mathrm{iid}}{\sim}}
\newcommand{\N}{\mathcal N}
\newcommand{\R}{\mathbb R}
\newcommand{\Prob}{\mathbb P}
\newcommand{\sol}{\par\textbf{حل.} }
\newcommand{\figcap}[1]{\par\smallskip{\small\color{muted}#1}\par}

\hypersetup{pdftitle={مرور سریع و راهنمای تطبیق},pdfauthor={دکتر سکینه دهقان؛ علیرضا جهانبخش و احسان شهبازی}}
\begin{document}
\appendix
\chapter{مرور سریع و راهنمای تطبیق}
\begin{center}{\small\color{muted} استاد درس: دکتر سکینه دهقان\quad | \quad گردآورندگان جزوه: علیرضا جهانبخش و احسان شهبازی}\end{center}

\begingroup\small\linespread{1.0}\selectfont
\section{انتخاب روش برای میانگین‌ها}
\begin{center}\small
\renewcommand{\arraystretch}{1.15}
\begin{tabular}{>{\raggedleft\arraybackslash}p{57mm} >{\centering\arraybackslash}p{59mm} >{\centering\arraybackslash}p{24mm}}
\toprule
مسئله & خطای استاندارد & توزیع مرجع\\
\midrule
یک میانگین؛ واریانس معلوم & $\sigma/\sqrt n$ & $Z$\\
یک میانگین نرمال؛ واریانس نامعلوم & $S/\sqrt n$ & $t_{n-1}$\\
اختلاف دو میانگین؛ واریانس‌ها معلوم & $\sqrt{\sigma_1^2/n_1+\sigma_2^2/n_2}$ & $Z$\\
دو نمونهٔ مستقل؛ واریانس نامعلوم مشترک & $S_p\sqrt{1/n_1+1/n_2}$ & $t_{n_1+n_2-2}$\\
دو نمونهٔ مستقل؛ ولچ & $\sqrt{S_1^2/n_1+S_2^2/n_2}$ & $t_\nu$ تقریبی\\
داده‌های زوجی & $S_D/\sqrt n$ & $t_{n-1}$\\
\bottomrule
\end{tabular}
\end{center}
برای ساخت بازهٔ دوطرفه، «برآورد نقطه‌ای $\pm$ چندک $1-\alpha/2$ ضرب‌در خطای استاندارد» را به کار ببرید. برای آزمون، «برآورد منهای مقدار فرض صفر، تقسیم بر خطای استاندارد» را با چندک مناسب جهت آزمون مقایسه کنید.
\section{چند تمایز که از خطا جلوگیری می‌کند}
\begin{important}{پیش از جای‌گذاری عددها}
\textbf{پارامتر دوم نرمال}: واریانس است؛ برای به دست آوردن $\sigma$ ریشه بگیرید.

\textbf{سطح و جهت آزمون}: دوطرفهٔ ۵ درصد به چندک $0.975$ نیاز دارد؛ یک‌طرفهٔ ۵ درصد به چندک $0.95$.

\textbf{استقلال و زوجی بودن}: نمونهٔ مستقل و زوج‌های مرتبط، خطاهای استاندارد متفاوتی دارند.

\textbf{واریانس نمونه}: $S^2$ با مخرج $n-1$، درحالی‌که برآورد حول میانگین معلوم با مخرج $n$ است.

\textbf{نسبت‌ها}: در آزمون تک‌نسبتی، $p_0$؛ در فاصلهٔ والد، $\widehat p$؛ در آزمون برابری دو نسبت، نسبت تجمیعی را به کار ببرید.

\textbf{نااریبی و سازگاری}: اولی دربارهٔ امید برآوردگر در حجم نمونهٔ معین است؛ دومی دربارهٔ رفتار دنبالهٔ برآوردگرها با رشد حجم نمونه.
\end{important}
\clearpage
\section{مهم‌ترین بازسازی‌ها و اصلاحات علمی}
موارد زیر تغییرهای مؤثر بر معنا یا نتیجه‌اند؛ بقیهٔ تغییرها عمدتاً یکدست‌سازی نمادگذاری، نگارش و تکمیل مراحل محاسبه هستند.
\begin{enumerate}\small\setlength{\itemsep}{2pt}\setlength{\parsep}{0pt}
\item \textbf{صفحه‌های ۱ و ۲:} روابط چگالی به توزیع‌های مطلقاً پیوسته محدود شد و قیود وجود گشتاورها افزوده شد.
\item \textbf{صفحهٔ ۴:} رابطهٔ درست بازگشتی تابع گاما، $\Gamma(\alpha+1)=\alpha\Gamma(\alpha)$، نوشته شد و قرارداد نرخ مشخص شد.
\item \textbf{صفحه‌های ۶ تا ۸:} اثبات اریبی واریانس با مخرج $n$، حل ترکیب بهینه و شرایط کافی سازگاری تکمیل شد. مثال دوجمله‌ای با آمارهٔ نامعمول، همراه با خوانش استفاده‌شده حفظ شد.
\item \textbf{صفحه‌های ۱۰ و ۱۲ تا ۱۳:} استقلال لازم برای ترکیب نرمال و ساخت کای‌دو و استیودنت تصریح شد؛ استدلال درجات آزادی و درجهٔ صحیح $n-1$ اصلاح شد.
\item \textbf{صفحهٔ ۱۴:} حجم نمونهٔ مورد نیاز به $62$ گرد شد؛ حاشیهٔ خطا به‌عنوان کران احتمالی، و نه تضمین قطعی، تعبیر شد.
\item \textbf{صفحهٔ ۱۶:} برای واریانس دوم، خوانش $60$ انتخاب و نتیجهٔ خوانش جایگزین $90$ نیز گزارش شد.
\item \textbf{صفحهٔ ۱۸:} روابط ولچ به‌عنوان تقریب معرفی شد.
\item \textbf{صفحهٔ ۲۲:} بررسی برابری میانگین‌ها از بررسی برابری واریانس‌ها تفکیک شد.
\item \textbf{صفحه‌های ۲۳ و ۳۹:} فرض نرمال بودن اختلاف‌های زوجی به‌درستی بیان شد.
\item \textbf{صفحهٔ ۲۴:} واریانس اختلاف‌ها $8.7$ و فاصلهٔ اطمینان آن $[-2.862,4.462]$ محاسبه شد؛ حذف جملهٔ $n\bar d^2$ در دست‌نویس اصلاح شد.
\item \textbf{صفحه‌های ۲۵ و ۲۶:} فاصله‌های نسبت به‌عنوان روش تقریبی معرفی و محدودیت نمونهٔ دارای پنج موفقیت مشخص شد.
\item \textbf{صفحهٔ ۳۱:} در مثال با آمارهٔ $2$، حجم نمونهٔ $25$ از مراحل حل بازسازی شد.
\item \textbf{صفحهٔ ۳۲:} آزمون دوطرفهٔ نمرات ریاضی به چندک $t_{0.975,8}=2.306$ اصلاح شد؛ تصمیم صحیح در سطح ۵ درصد، رد نکردن فرض صفر است.
\item \textbf{صفحه‌های ۳۳ و ۳۴:} جهت تمام اختلاف میانگین‌ها به «اول منهای دوم» یکدست شد؛ جهت نامساوی‌ها نیز متناظر تنظیم شد.
\item \textbf{صفحه‌های ۴۳ تا ۴۵:} مقدار دقیق‌تر $p=0.0227501$ نوشته شد و تصمیم‌های سطح ۵ و ۱ درصد از هم تفکیک شدند. تفسیر «رد نکردن» جایگزین «اثبات یا قبول قطعی» شد.
\end{enumerate}
\clearpage
\section{نقشهٔ تطبیق با فایل ۴۵ صفحه‌ای}
شماره‌های ستون اول، شمارهٔ صفحات PDF ارسالی هستند؛ ممکن است با شماره‌های دست‌نویس گوشهٔ برگه‌ها یکسان نباشند. این جدول امکان بازبینی پوشش همهٔ صفحات را فراهم می‌کند.
\begin{center}\small
\renewcommand{\arraystretch}{1.2}
\begin{longtable}{>{\centering\arraybackslash}p{22mm} p{113mm}}
\toprule
صفحات مبنا & محل و موضوع بازنویسی\\
\midrule\endfirsthead
\toprule صفحات مبنا & محل و موضوع بازنویسی\\\midrule\endhead
۱–۲ & فصل اول: تعریف‌ها، تابع احتمال و چگالی، توزیع مشترک و امید ریاضی\\
۳–۴ & فصل اول: واریانس، کوواریانس، میانگین نمونه و توزیع‌های مهم\\
۵ & فصل دوم: جامعه، نمونه، آماره و مثال توزیع نمونه‌ای میانگین\\
۶ & فصل دوم: نااریبی، واریانس نمونه، خطای میانگین مربعات و کارایی\\
۷–۸ & فصل دوم: مثال‌های کارایی، ترکیب خطی و سازگاری؛ ارجاع‌های کتاب\\
۹–۱۱ & فصل سوم: فاصلهٔ اطمینان، کمیت محوری و میانگین با واریانس معلوم\\
۱۲–۱۳ & فصل سوم: کای‌دو و استیودنت؛ میانگین با واریانس نامعلوم\\
۱۴ & فصل سوم: حاشیهٔ خطا و تعیین حجم نمونه\\
۱۵ & فصل سوم: قضیهٔ حد مرکزی و اختلاف دو میانگین\\
۱۶–۱۸ & فصل سوم: اختلاف میانگین با واریانس‌های معلوم، مشترک و ولچ\\
۱۹–۲۰ & فصل سوم: فاصلهٔ واریانس با میانگین معلوم و نامعلوم\\
۲۰–۲۲ & فصل سوم: توزیع $F$، نسبت واریانس‌ها و خاصیت معکوس چندک\\
۲۳–۲۴ & فصل سوم: داده‌های زوجی و مثال نمرهٔ دو درس\\
۲۴–۲۵ & فصل سوم: فاصلهٔ اطمینان یک نسبت و مثال قطعات معیوب\\
۲۶ & فصل سوم: اختلاف دو نسبت و مثال نظر دانشجویان\\
۲۷–۲۸ & فصل چهارم: فرض‌ها، خطاها، توان و مثال سکه\\
۲۹–۳۱ & فصل چهارم: سه جهت آزمون میانگین با واریانس معلوم\\
۳۲ & فصل چهارم: آزمون استیودنت تک‌نمونه‌ای و مثال نمرات ریاضی\\
۳۳–۳۴ & فصل چهارم: آزمون اختلاف میانگین‌های مستقل\\
۳۵–۳۶ & فصل چهارم: آزمون واریانس یک جامعهٔ نرمال\\
۳۷–۳۸ & فصل چهارم: آزمون برابری واریانس دو جامعهٔ نرمال\\
۳۹ & فصل چهارم: آزمون میانگین در داده‌های زوجی\\
۴۰ & فصل چهارم: آزمون یک نسبت\\
۴۱–۴۲ & فصل چهارم: آزمون برابری دو نسبت و برآورد تجمیعی\\
۴۳–۴۵ & فصل چهارم: مقدار احتمالی، مثال محصول جدید و تصمیم در دو سطح\\
\bottomrule
\end{longtable}
\end{center}
\endgroup

\end{document}
